\section{Buggen ett enhetstest inte kan se}
\label{sec:interaction}

Ett \textbf{enhetstest} testar en enhet isolerad. Ett \textbf{komponenttest} testar flera delar som
samverkar, fortfarande utan hårdvara: här \code{app::EchoNode} plus \code{driver::can::Stub},
kopplade genom \code{driver::can::Interface}.

Skillnaden syns tydligast i vilken sorts bugg de fångar. Anta att \code{run()} skrivs så här:

\begin{cppcode}
Frame frame{};
if (!myDriver.receive(frame)) { return; }

Frame response{};
response.id  = myResponseId;
response.dlc = myDriver.receive(frame) ? frame.dlc : 0U;   // fel: anropar receive() igen
\end{cppcode}

Varje enhetstest är fortfarande grönt. \code{Stub} gör exakt det den lovar, och \code{Frame}
likaså. Buggen ligger i \textbf{interaktionen}: det andra anropet konsumerar tillståndet det första
lämnade.

\textbf{Det är hela poängen.} Interaktionsbuggar --- ett anrop för mycket, ett anrop i fel ordning,
ett returvärde som ignoreras --- är den vanligaste sortens fel i drivrutinskod, och de bor mellan
enheterna, inte i dem.

\section{\texorpdfstring{\code{app::EchoNode}}{app::EchoNode}}
\label{sec:echonode}

En applikationskomponent som gör en sak: läser en frame och ekar tillbaka den med ett annat ID.

\begin{cppcode}
namespace driver::can { class Interface; }

namespace app
{
class EchoNode
{
public:
    explicit EchoNode(driver::can::Interface& driver, std::uint16_t responseId) noexcept;

    void run() noexcept;
    [[nodiscard]] std::uint32_t echoCount() const noexcept;

private:
    driver::can::Interface& myDriver; /**< Driver used for all traffic. Not owned. */
    std::uint16_t myResponseId;
    std::uint32_t myEchoCount;
};
} // namespace app
\end{cppcode}

\subsection{Lagergränsen, i en rad}

Lägg märke till vad headern \emph{inte} inkluderar: \code{stub.hpp}, \code{spi.hpp},
\code{byte_transport.hpp}. Bara en framåtdeklaration.

\textbf{Den raden är arkitekturens gräns uttryckt i kod.} En granskare som ser
\code{#include "driver/can/spi.hpp"} här ska stoppa PR:en direkt. Regeln, som gäller hela projektet:
ingen fil ovanför \code{driver/can/interface.hpp} får innehålla ordet SPI.

\subsection{Dependency injection}

\code{EchoNode} \textbf{skapar inte} sin drivrutin. Den får den. Alternativet är värt att titta på
för att se vad det kostar:

\begin{cppcode}
class EchoNode
{
public:
    EchoNode() : myDriver{} {}   // skapar sin egen Spi
private:
    driver::can::Spi myDriver;   // nej
};
\end{cppcode}

Den varianten är omöjlig att testa: det finns inget sätt att ge den en stubb, alltså inget sätt
att köra den utan ett FPGA på skrivbordet. Den binder dessutom applikationen till en konkret
drivrutinstyp, vilket gör hela \code{Interface} meningslös.

\textbf{Varför en referens och inte en pekare?} En referens kan inte vara null och kan inte bytas
efter konstruktion. Båda är egenskaper vi vill ha.

\textbf{Varför inte \code{std::unique_ptr}?} Därför att \code{EchoNode} inte \emph{äger}
drivrutinen. Ägandet ligger hos factoryn i \kapref{ch:avr}, och drivrutinens livstid är längre än
komponentens. Doxygen-kommentaren säger uttryckligen "Not owned", och det är värt att skriva ut.

\subsection{Två detaljer i \texorpdfstring{\code{run()}}{run()}}

\textbf{Loopen går till \code{frame.dlc}, inte till \code{MaxDataLen}.} Att kopiera åtta bytes hade
"fungerat" eftersom resten ändå är nollor --- men \code{response.dlc} säger hur många som gäller,
och kod som kopierar mer än den säger sig kopiera slutar stämma när någon annan förutsättning
ändras. Exakt samma resonemang som maskeringen i \kapref{ch:driver}.

\textbf{\code{myEchoCount} räknas upp bara när \code{send()} lyckades.} Räknaren betyder "ekade
frames", inte "försök".

\section{Att läsa komponenttestet}
\label{sec:readcomp}

\begin{cppcode}
TEST(EchoNode, echoesInjectedFrameWithResponseId)
{
    // Arrange.
    driver::can::Stub stub{};
    app::EchoNode node{stub, 0x200U};

    driver::can::Frame incoming{};
    incoming.id      = 0x123U;
    incoming.dlc     = 2U;
    incoming.data[0] = 0xAAU;
    incoming.data[1] = 0xBBU;
    stub.inject(incoming);

    // Act.
    node.run();

    // Assert.
    driver::can::Frame echoed{};
    EXPECT_TRUE(stub.receive(echoed));
    EXPECT_EQ(echoed.id, 0x200U);
    EXPECT_EQ(node.echoCount(), 1U);
}
\end{cppcode}

Testet har två roller samtidigt: det \textbf{matar in} trafik med \code{inject()}, och det
\textbf{inspekterar} utgående trafik med \code{receive()} --- båda genom samma testdubblare.

Vad den utdelade sviten täcker:

\begin{booktable}{@{}lL@{}}
\thead{Fall} & \thead{Vad det bevisar} \\ \midrule
Ingen injicerad frame & \code{echoCount()} oförändrad. \\
En injicerad frame & Rätt \code{id}, \code{dlc} och data ekas; räknaren går upp med ett. \\
Två frames, två \code{run()} & Räknaren blir 2. \\
En frame, \code{run()} \textbf{två} gånger & Räknaren blir 1 --- inte 2. Fångar buggen i
  \secref{sec:interaction}. \\
\code{dlc = 0} och \code{dlc = 8} & Gränsvärdespar. \\
\code{send()} misslyckas & Räknaren går inte upp. \\
\end{booktable}

\subsection{Samma mönster, ett lager ned}

Det här är det enda du behöver bära med dig härifrån. I \kapref{ch:testning} gör den utdelade
sviten \textbf{exakt samma sak}, en söm längre ned: \code{BankTransport} ersätter \code{Stub},
\code{Spi} är det som testas, och testet skriptar vad hårdvaran svarar i stället för vilka frames
som kommer in.

Att det \emph{är} samma mönster är ingen tillfällighet --- det är vad man får av att lägga sömmarna
på rätt ställen.

\subsection{Vad det inte bevisar}

Att \code{EchoNode} samverkar rätt med \emph{någon} implementation av \code{Interface}. Ingenting
om \code{Spi}, som ännu inte finns, och ingenting om huruvida registerkartan lästs rätt.
